\SourceLeaf{150}Vorteile vereint erhält. Die\,Methode\,der\,kleinsten\ Quadrate erweist sich in allen Fällen und bei jedem Fehlergesetz als die absolut vorteilhafteste. Der Vergleich der Genauigkeit der Ergebnisse mit der der Beobachtungen, den ich in meiner Theoria auf das Fehlergesetz $e^{-hhxx}$ gegründet hatte, bleibt allgemein gültig. Zugleich hat man den Vorteil, dass sich alles durch sehr klare und einfache analytische Entwicklungen beweisen und lösen lässt. Das scheint beim Prinzip und bei der Behandlung von \textsc{Laplace} keineswegs der Fall zu sein. Insbesondere scheint die Verallgemeinerung seines Schlusses von 2 unbekannten Größen auf eine beliebige Anzahl noch nicht hinreichend einleuchtend zu sein.

\noindent . . . . . . . . . . .



\Needspace{6\baselineskip}
\begin{center}
\textsc{Gauss} an \textsc{Schumacher}. Göttingen, 2. Februar 1825.
\end{center}

\noindent . . . . . . . . . . Die sogenannten wahrscheinlichen Fehler würde ich eigentlich ganz verwerfen, da sie von einer Hypothese abhängen. Man kann sie jedoch berechnen, indem man die mittleren Fehler mit 0,6744897 multipliziert.

\noindent . . . . . . . . . . .



\Needspace{6\baselineskip}
\begin{center}
\textsc{Gauss} an \textsc{Olbers}. Göttingen, 15. März 1827.
\end{center}

\noindent Ich bin Ihnen sehr verbunden, mein liebster Freund, dass Sie mir die drei Hefte des Philosophical Magazine freundlicherweise geschickt haben. Ich habe die Aufsätze von \textsc{Ivory} durchgesehen und müsste sagen, dass sie mich sehr befremden, wenn ich nicht bereits durch die Abhandlung in den Philosophical Transactions (1824), die ich in meinem letzten Brief erwähnte, darauf vorbereitet gewesen wäre. Ich habe Herrn I.\ seit Langem als scharfsinnigen Mathematiker mit großem Geschick im Rechnen geschätzt. Von \textls{diesem} Urteil gehe ich auch jetzt nicht ab. Im Hauptpunkt der erwähnten Abhandlung vermisse ich allerdings den logischen\space%
